\newpage
\chapter{特征值}
\section{特征值与特征向量}
\subsection{引入}
\begin{red}
\begin{remark}
    考虑一个曾提出的问题,如何找到一个线性空间的一组基,使得该线性空间上所有线性变换的表示矩阵尽量简单（比如对角阵）.或者说,能否找到一些尽量简单的矩阵使得其他矩阵均与它们中的某一个相似（比如对角阵）.
\end{remark}
\end{red}
设$\mathbb{K}$上的$n$维线性空间$V$,$\bm{\varphi}\in\mathcal{L}\left(V\right)$,且$\bm{\varphi}$在$\left\{
    \bm{e}_1,\bm{e}_2,\cdots,\bm{e}_n
\right\}$这组基下的表示矩阵为$\bm{A}=\mathrm{diag}\,\left\{
    \lambda_1,\lambda_2,\cdots,\lambda_n
\right\}$,从而
\[
\left(
       \bm{\varphi}\left(\bm{e}_1\right),\bm{\varphi}\left(\bm{e}_2\right),\cdots,\bm{\varphi}\left(\bm{e}_n\right)   
\right)=
\left(
    \bm{e}_1,\bm{e}_2,\cdots,\bm{e}_n
\right)\begin{pmatrix}
    \lambda_1&0&\cdots&0\\
    0&\lambda_2&\cdots&0\\
    \vdots&\vdots&\ddots&\vdots\\
    0&0&\cdots&\lambda_n
\end{pmatrix}
\]即$\bm{\varphi}\left(
    \bm{e}_i
\right)=\lambda_i\bm{e}_i,\forall 1\leqslant i\leqslant n.$那么
\[
\bm{\alpha}= \sum_{i=1}^{n} a_i\bm{e}_i\Longrightarrow
\bm{\varphi}\left(
    \bm{\alpha}
\right)=\sum_{i=1}^{n}a_i\lambda_i
    \bm{e}_i
\]

设$\lambda_1,\cdots,\lambda_r$非零,$
\lambda_{r+1}=\cdots=\lambda_n=0.$则$\mathrm{r}\left(
    \bm{\varphi}
\right)=r$且$\mathrm{Ker}\,\bm{\varphi}=L\left(
     \bm{e}_{r+1},\cdots,\bm{e}_n
\right),\mathrm{Im}\,\bm{\varphi}=L\left(
    \bm{e}_1,\cdots,\bm{e}_r
\right)$.

这一切的关键在于$
\bm{\varphi}\left(
    \bm{e}_i
\right)=\lambda_i\bm{e}_i,\bm{e}_i\neq 0,\forall 1\leqslant i\leqslant n
$.这样的$\lambda_i$称为线性变换$\bm{\varphi}$的特征值,$\bm{e}_i$称为$\bm{\varphi}$的特征向量.
\subsection{特征值与特征向量}
\begin{formal}
    \begin{definition}[线性变换意义的特征值与特征向量]\label{def:线性变换意义的特征值与特征向量}
        设$\bm{\varphi}\in\mathcal{L}\left(
            V_{\mathbb{K}}^n
        \right)$,若存在$\lambda\in\mathbb{K},\bm{0}\neq \bm{e}\in V\,\textup{s.t.}$
        \[
        \bm{\varphi}\left(\bm{e}\right)=\lambda\bm{e}
        \]则称$\lambda$为$\bm{\varphi}$的特征值,$\bm{e}$为$\bm{\varphi}$对应于
        特征值$\lambda$的特征向量.
    \end{definition}
\end{formal}
\begin{formal}
\begin{definition}[线性变换意义的特征子空间]\label{def:线性变换意义的特征子空间}
    对于每一个特征值$\lambda$,其对应的特征向量加上一个零向量构成的集合称为$\bm{\varphi}$关于特征值$\lambda$的特征子空间,记作$V_{\lambda}\left(\bm{\varphi}\right)=
        \left\{
            \bm{v}\in V\mid \bm{\varphi}\left(\bm{v}\right)=\lambda\bm{v}
        \right\}$.
\end{definition}
\end{formal}
\begin{formal}
\begin{theorem}[特征子空间的维数]\label{thm:特征子空间的维数}
    根据\cref{def:线性变换意义的特征子空间}知特征子空间$V_{\lambda_0}$的维数$\mathrm{dim}\,V_{\lambda_0}$等于$\bm{\varphi}$的特征值$\lambda_0$的几何重数.
\end{theorem}
\end{formal}
\begin{green}
\begin{lemma}[特征子空间]\label{lem:特征子空间}
    特征子空间$V_{\lambda}\left(\bm{\varphi}\right)$是$V$的子空间,并且是不变子空间.
\end{lemma}
\end{green}
下面考虑代数的角度.
一般地,设$\bm{\varphi}$在某一组基下的表示矩阵为$\bm{A}\in M_n\left(
    \mathbb{K}
\right)$,$\bm{e}$的坐标向量为$\bm{\alpha}\in\mathbb{K}^n$,则$\bm{\varphi}\left(
    \bm{e}
\right)$的坐标向量为$\bm{A\alpha}$,于是
\[
\bm{\varphi}\left(\bm{e}\right)=\lambda\bm{e}\Longleftrightarrow
\bm{A\alpha}=\lambda\bm{\alpha}
\]
\begin{formal}
\begin{definition}[矩阵意义的特征值与特征向量]\label{def:矩阵意义的特征值与特征向量}
    设$\bm{A}\in M_n\left(
        \mathbb{K}
    \right)$,若存在$\lambda\in\mathbb{K},\bm{0}\neq\bm{\alpha}\in\mathbb{K}^n\,\textup{s.t.}$
    \[
    \bm{A\alpha}=\lambda\bm{\alpha}
    \]则称$\lambda$为$\bm{A}$的特征值,$\bm{\alpha}$为$\bm{A}$对应于特征值$\lambda$的特征向量.
\end{definition}
\end{formal}
\begin{formal}
\begin{definition}[矩阵意义的特征子空间]\label{def:矩阵意义的特征子空间}
    对于每一个特征值$\lambda$,定义$V_{\lambda}$为线性方程组
    \[
    \left(
        \lambda\bm{I}_n-\bm{A}
    \right)\bm{\alpha}=\bm{0}
    \]
    的解空间,称为$\bm{A}$关于特征值$\lambda$的特征子空间.
\end{definition}
\end{formal}
\begin{red}
\begin{remark}
    $\lambda_0$为$\bm{A}$的特征值当且仅当
    $\exists \bm{0}\neq \bm{\alpha}\in\mathbb{K}^n \,\textup{s.t.}\bm{A\alpha}=\lambda_0\bm{\alpha}$当且仅当
    线性方程组$\left(
        \lambda_0\bm{I}_n-\bm{A}
    \right)\bm{\alpha}=\bm{0}$有非零解当且仅当$\left(
        \lambda_0\bm{I}_n-\bm{A}
    \right)$非奇异当且仅当
    $\left|\lambda_0\bm{I}_n-\bm{A}\right|=0$.

    即$\lambda$是方程
    \[\left|\lambda\bm{I}_n-\bm{A}\right|=0\]的根.即考虑
    \[
    \left|
         \lambda\bm{I}_ n-\bm{A}
    \right|=\begin{vmatrix}
        \lambda-a_{11}&-a_{12}&\cdots&-a_{1n}\\
        -a_{21}&\lambda-a_{22}&\cdots&-a_{2n}\\
        \vdots&\vdots&\ddots&\vdots\\
        -a_{n1}&-a_{n2}&\cdots&\lambda-a_{nn}
    \end{vmatrix}
    \]
    这是一个关于$\lambda$的$n$次首一多项式,称为$\bm{A}$的特征多项式.
\end{remark}
\end{red}
\begin{formal}
\begin{definition}[矩阵的特征多项式]\label{def:矩阵的特征多项式}
    设$\bm{A}\in M_n\left(\mathbb{K}\right)$,称$\left|
        \lambda\bm{I}_n-\bm{A}
    \right|$为$\bm{A}$的特征多项式.
\end{definition}
\end{formal}
\begin{green}
\begin{lemma}[相似矩阵的特征根]\label{lem:相似矩阵的特征根}
    相似矩阵有相同的特征多项式,从而有相同的特征值（计重数）.
\end{lemma}
\begin{Proof}
设$\bm{B}=\bm{P}^{-1}\bm{AP}$,则
\begin{align*}
    \left|
       \lambda\bm{I}_n-\bm{B}   
\right|&=
\left|
    \lambda\bm{I}_n-\bm{P}^{-1}\bm{AP}
\right|=\left|
    \bm{P}^{-1}\left(
        \lambda\bm{I}_n-\bm{A}
    \right)\bm{P}
\right|\\
&=\left|
    \bm{P}^{-1}
\right|\left|
    \lambda\bm{I}_n-\bm{A}
\right|\left|
    \bm{P}
\right|=\left|
    \lambda\bm{I}_n-\bm{A}
\right|\qedhere
\end{align*}
\end{Proof}
\end{green}
\begin{formal}
\begin{definition}[线性变换的特征多项式]\label{def:线性变换的特征多项式}
    设$\bm{\varphi}\in\mathcal{L}\left(
        V_{\mathbb{K}}^n
    \right)$,任取一组基下的表示矩阵$\bm{A}$,
    定义$\bm{\varphi}$的特征多项式是$\bm{A}$的特征多项式（由上知特征多项式的定义不依赖于基或表示矩阵的选取）.记作$\left|
        \lambda\bm{I}_V-\bm{\varphi}
    \right|$.
\end{definition}
\end{formal}
\begin{green}
\begin{lemma}[特征值的和与积]\label{lem:特征值的和与积}
    设$\bm{A}$的特征值为$\lambda_1,\lambda_2,\cdots,\lambda_n$,则
    \begin{align*}
        \lambda_1+\lambda_2+\cdots+\lambda_n&=\mathrm{tr}\,\bm{A}\\
        \lambda_1\lambda_2\cdots\lambda_n&=\left|\bm{A}\right|
    \end{align*}
\end{lemma}
\begin{Proof}
设\begin{align*}
    \left|\lambda\bm{I}_n-\bm{A}\right|&=
    \lambda^n+a_1\lambda^{n-1}+\cdots+a_{n-1}\lambda+a_n\\
    &=\left(
        \lambda-\lambda_1
    \right)\left(
        \lambda-\lambda_2
    \right)\cdots\left(
        \lambda-\lambda_n
    \right)
\end{align*}
其中$a_1=
-\left(
    a_{11}+a_{22}+\cdots+a_{nn}
\right)=-\mathrm{tr}\,\bm{A}$,令$\lambda = 0$得$
a_n=\left(-1\right)^n\left|\bm{A}\right|
$.考虑\cref{thm:Vieta定理}Veita定理即可.
\end{Proof}
\end{green}
\begin{green}
\begin{corollary}[$r$阶主子式之和]\label{cor:r阶主子式之和}
    进一步,可以证明
\begin{align*}
    \sum_{1\leqslant i_1<i_2<\cdots<i_r\leqslant n}\lambda_{i_1}\lambda_{i_2}\cdots\lambda_{i_r}=
    \sum_{1\leqslant i_1<i_2<\cdots<i_r\leqslant n} \bm{A}\begin{pmatrix}
        i_1&i_2&\cdots&i_r\\
        i_1&i_2&\cdots&i_r
    \end{pmatrix}
\end{align*}
\end{corollary}
\end{green}
\begin{green}
\begin{corollary}[非异阵判定]\label{cor:非异阵判定}
    设$\bm{A}\in M_n\left(\mathbb{K}\right)$,则$\bm{A}$非异当且仅当
    $\lambda_i\neq 0,\forall 1\leqslant i\leqslant n.$
\end{corollary}
\begin{Proof}
\[
 \lambda_1\lambda_2\cdots \lambda_n=\left|\bm{A}\right|\qedhere
\]
\end{Proof}
\end{green}
\begin{formal}
\begin{proposition}[求特征值和特征向量的方法]\label{prop:求特征值和特征向量的方法}\,

     \begin{enumerate} [label = {\textup{(\arabic*)}}]
        \item 写出特征矩阵$\lambda \bm{I}_n-\bm{A}$,求出特征多项式$\left|\lambda \bm{I}_n-\bm{A}\right|=0$的根$\lambda_1,\lambda_2,\cdots,\lambda_n$即为特征值.
        \item $\forall 1\leqslant i\leqslant n$,解线性方程组$\left(
            \lambda_i\bm{I}_n-\bm{A}
        \right)\bm{\alpha}=\bm{0}$,得到的非零解即为对应于$\lambda_i$的特征向量.
     \end{enumerate}
\end{proposition}
\end{formal}
\begin{red}
\begin{remark}
    对于$\bm{A}\in M_n\left(
        \mathbb{K}
    \right)$,其特征值可能不在$\mathbb{K}$中,例如矩阵
    $\begin{bmatrix}
        0&-1\\
        1&0
    \end{bmatrix}\in M_n\left(
        \mathbb{R}
    \right)$的特征值为$\pm \mathrm{i}
    \notin \mathbb{R}$.因此更一般地,我们考虑特征值相关问题时通常在复数域$\mathbb{C}$上考虑.
\end{remark}
\end{red}
\begin{brown}
\begin{example}
    考虑上三角阵
    \[
    \bm{A}=\begin{pmatrix}
        a_{11}&a_{12}&\cdots&a_{1n}\\
        0&a_{22}&\cdots&a_{2n}\\
        \vdots&\vdots&\ddots&\vdots\\
        0&0&\cdots&a_{nn}
    \end{pmatrix}
    \]其特征多项式显然为
    \[
    \left|
    \lambda\bm{I}_n-\bm{A}    
    \right|=
    \begin{vmatrix}
        \lambda-a_{11}&-a_{12}&\cdots&-a_{1n}\\
        0&\lambda-a_{22}&\cdots&-a_{2n}\\
        \vdots&\vdots&\ddots&\vdots\\
        0&0&\cdots&\lambda-a_{nn}
    \end{vmatrix}
    =\left(
        \lambda-a_{11}
    \right)\left(
        \lambda-a_{22}
    \right)\cdots\left(
        \lambda-a_{nn}
    \right)
    \]特征值为
    $\lambda_1=a_{11},\lambda_2=a_{22},\cdots,\lambda_n=a_{nn}$.
\end{example}
\end{brown}
\begin{brown}
\begin{example}
    考虑矩阵
    \[
    \begin{bmatrix}
        3&1&-1\\
        2&2&-1\\
        2&2&0
    \end{bmatrix}
    \]
    则其特征根
    $\lambda_1=1,\lambda_2=\lambda_3=2$.求解
    \[
    \left(
           \lambda\bm{I}_3-\bm{A}   
    \right)\bm{x}=\bm{0}
    \]得
    ：$\lambda_1=1$时
    \[
      \begin{bmatrix}
        -2&-1& 1\\
        -2&-1&1\\
        -2&-2&1
      \end{bmatrix}\longrightarrow
      \begin{bmatrix}
        1&0&-\frac{1}{2}\\
        0&1&0\\
        0&0&0
      \end{bmatrix}\Longrightarrow
      \bm{\xi}_1=\begin{pmatrix}
        1\\0\\2
      \end{pmatrix}
    \]
    对应特征向量为$c_1\bm{\xi}_1,c\neq 0.$
    $\lambda_2=\lambda_3=2$时
    \[
    \begin{bmatrix}
        -1&-1&1\\
        -2&0&1\\
        -2&-2&2
    \end{bmatrix}\longrightarrow
    \begin{bmatrix}
        1&0&-\frac{1}{2}\\
        0&1&-\frac{1}{2}\\
        0&0&0
    \end{bmatrix}\Longrightarrow
    \bm{\xi}_2=\begin{pmatrix}
        1\\1\\2
    \end{pmatrix}
    \]对应特征向量为$c_2\bm{\xi}_2,c_2\neq 0$.
\end{example}
\end{brown}
\subsection{相似}
\begin{formal}
\begin{theorem}[必定复相似]\label{thm:必定复相似}
    任一复方阵必复相似于一个上（下）三角阵.
\end{theorem}
\begin{Proof}
设$\bm{A}\in M_n\left(
    \mathbb{C}
\right)$,对阶数$n$进行归纳,当$n=1$时结论显然成立.假设$n-1$时结论成立,考虍$n$时的情况.由\cref{thm:代数学基本定理}代数学基本定理,任取$\bm{A}$的特征值$\lambda_1$和特征向量$\bm{0}\neq \bm{\alpha}_1\in \mathbb{C}^n$即
\[
\bm{A\alpha}_1=\lambda_1\bm{\alpha}_1
\]

根据\cref{thm:基扩张定理}基扩张定理,将$\bm{\alpha}_1$扩张为全空间$\mathbb{C}^n$的一组基$\left\{
    \bm{\alpha}_1,\bm{\alpha}_2,\cdots,\bm{\alpha}_n
\right\}$.令$\bm{P}=\left(
    \bm{\alpha}_1,\bm{\alpha}_2,\cdots,\bm{\alpha}_n
\right)\in M_n\left(
    \mathbb{C}
\right)$且$\bm{P}$非异.因为
\begin{align*}
    \bm{AP}&=\left(
        \bm{A\alpha}_1,\bm{A\alpha}_2,\cdots,\bm{A\alpha}_n
    \right)\\
    &=\left(
        \bm{\alpha}_1,  \bm{A\alpha}_2,\cdots,\bm{A\alpha}_n
    \right)\begin{bmatrix}
        \lambda_1&*\\
        \bm{O}&\bm{A}_{n-1}
    \end{bmatrix}
\end{align*}
其中$\bm{A}_{n-1}\in M_{n-1}\left(
    \mathbb{C}
    \right)$.于是
    \[
    \bm{P}^{-1}\bm{AP}=\begin{pmatrix}
        \lambda_1&*\\
        \bm{O}&\bm{A}_{n-1}
    \end{pmatrix}
    \]
    由归纳假设,$\exists\textup{非异阵} 
    \bm{Q}\in M_{n-1}\left(
        \mathbb{C}
    \right)\,\textup{s.t.}$
    \[
    \bm{Q}^{-1}\bm{A}_{n-1}\bm{Q}=\begin{pmatrix}
        
        \lambda_2&\cdots&*\\
        \vdots&\ddots&\vdots\\
        \bm{O}&\cdots&\lambda_{n}
    \end{pmatrix}
    \]作非异阵$
    \bm{R}=\begin{pmatrix}
        1&\bm{O}\\
        \bm{O}&\bm{Q}
    \end{pmatrix}\in M_n\left(
        \mathbb{C}
    \right)$,则
    \begin{align*}
        \bm{R}^{-1}\bm{P}^{-1}\bm{APR}&=
        \bm{R}^{-1}\left(
            \bm{P}^{-1}\bm{A}\bm{P}
        \right)\bm{R}\\
        &=\begin{bmatrix}
            1&\bm{O}\\
            \bm{O}&\bm{Q}^{-1}
        \end{bmatrix}\begin{bmatrix}
            \lambda_1&*\\
            \bm{O}&\bm{A}_{n-1}
        \end{bmatrix}\begin{bmatrix}
            1&\bm{O}\\
            \bm{O}&\bm{Q}
        \end{bmatrix}\\
        &=
        \begin{bmatrix}
            \lambda_1&\cdots&*\\
            \vdots&\ddots&\vdots\\
            \bm{O}&\cdots&\lambda_n
        \end{bmatrix}\qedhere
    \end{align*}
\end{Proof}
\end{formal}
\begin{green}
\begin{corollary}[限制在某一个数域]\label{cor:限制在某一个数域}
    设$\bm{A}\in M_n\left(
        \mathbb{K}
    \right)$的特征值均在$\mathbb{K}$中,则$\exists 
    \textup{非异阵}\bm{P}\in M_n\left(
        \mathbb{K}
    \right)\,\textup{s.t.}$
    \[
    \bm{P}^{-1}\bm{AP}=\begin{pmatrix}
        \lambda_1&*&\cdots&*\\
        \bm{O}&\lambda_2&\cdots&*\\
        \vdots&\vdots&&\vdots\\
        \bm{O}&\bm{O}&\cdots&\lambda_n
    \end{pmatrix}
    \]
\end{corollary}
\begin{Proof}
证明基本同\cref{thm:必定复相似}.
\end{Proof}
\end{green}
\subsection{相似关系}
\begin{red}
\begin{remark}
    一个问题是,如果知道了$\bm{A}\in M_n\left(
        \mathbb{K}
    \right)$的特征值,那么
    \[
    f\left(
        \bm{A}
    \right)=a_m\bm{A}^m+a_{m-1}\bm{A}^{m-1}+\cdots+a_1\bm{A}+a_0\bm{I}_n\in M_n\left(
        \mathbb{K}
    \right)
    \]的特征值是否也能够知道.
\end{remark}
\end{red}
\begin{green}
\begin{lemma}[幂保持]\label{lem:幂保持}
    相似关系在幂运算下保持.
    \[
    \left(
        \bm{P}^{-1}\bm{AP}   
    \right)^m=\bm{P}^{-1}\bm{A}^m\bm{P}
    \]
\end{lemma}
\end{green}
\begin{green}
\begin{corollary}[相似与多项式相容]\label{cor:相似与多项式相容}
    由\cref{lem:幂保持},有
    \[
f\left(
    \bm{P}^  {-1}\bm{AP}
\right)= \bm{P}^{-1}f\left(
    \bm{A}
\right)\bm{P}
\]即相似与多项式相容.
\end{corollary}
\end{green}
\begin{green}
\begin{corollary}[求逆保持]\label{cor:求逆保持}
    相似关系在求逆运算下保持.
    \[
    \left(
        \bm{P}^{-1}\bm{AP}
    \right)^{-1}=\bm{P}^{-1}\bm{A}^{-1}\bm{P}
    \]
\end{corollary}
\end{green}
\begin{green}
\begin{corollary}[求伴随保持]\label{cor:求伴随保持}
    相似关系在求伴随运算下保持.
    \[
    \left(
        \bm{P}^{-1}\bm{AP}
    \right)^*=\bm{P}^*\bm{A}^*\left(
        \bm{P}^*
    \right)^{-1}
    \]
\end{corollary}
\end{green}
\subsection{运算的特征值}
\begin{formal}
\begin{proposition}[特征值保持]\label{prop:特征值保持}
    设$\bm{A}\in M_n\left(
        \mathbb{K}
    \right)$的特征值为$\lambda_1,\lambda_2,\cdots,\lambda_n$,则$f\left(
        \bm{A}
    \right)$的特征值为$f\left(
        \lambda_1
    \right),f\left(
        \lambda_2
    \right),\cdots,f\left(
        \lambda_n
    \right)$.
\end{proposition}
\begin{Proof}
由前文,存在非异阵$\bm{P}\in M_n\left(
    \mathbb{K}
\right)$使得
\[
\bm{P}^{-1}\bm{AP}=\begin{pmatrix}
    \lambda_1&\cdots&*\\
    \vdots&\ddots&\vdots\\
    \bm{O}&\cdots&\lambda_n
\end{pmatrix}
\]考虑
$
\bm{P}^{-1}f\left(
    \bm{A}
\right)\bm{P}=f\left(
    \bm{P}^{-1}\bm{AP}
\right)
$,因为上三角阵的加减乘运算仍为上三角阵,且运算的主对角线元素等于主对角线元素的运算,则
\[
f\left(
    \bm{P}^{-1}\bm{AP}
\right)=\begin{pmatrix}
    f\left(
        \lambda_1
    \right)&\cdots&*\\
    \vdots&\ddots&\vdots\\
    \bm{O}&\cdots&f\left(
        \lambda_n
    \right)
\end{pmatrix}\qedhere
\]
\end{Proof}
\end{formal}
\begin{formal}
\begin{proposition}[多项式保持]\label{prop:多项式保持}
    设$\bm{A}\in M_n\left(
        \mathbb{K}
    \right)$适合$g\left(
        x
    \right)\in \mathbb{K}\left[x\right]$即$g\left(\bm{A}\right)=\bm{O}$,则$\bm{A}$的任意特征值也适合该多项式即$g\left(
        \lambda_i
    \right)=0,\forall 1\leqslant i\leqslant n$.
\end{proposition}
\begin{Proof}
设$\bm{A\alpha}=\lambda\bm{\alpha},\bm{0}\neq \bm{\alpha}\in \mathbb{C}^n$,则$\bm{A}^m\bm{\alpha}=\lambda^m\bm{\alpha}$.于是
\[
\bm{O}=g\left(\bm{A}\right)\bm{\alpha}=g\left(\lambda\right)\bm{\alpha}\Longrightarrow
g\left(\lambda\right)=0
\]

亦可用\cref{prop:特征值保持}完成,因为$\bm{O}=g\left(\bm{A}\right)$的特征值为$g\left(\lambda_1\right),
g\left(\lambda_2\right),\cdots,g\left(\lambda_n\right)
$且零矩阵的特征值均为零,得证.
\end{Proof}
\end{formal}
\begin{formal}
\begin{proposition}[逆阵的特征值]\label{prop:逆阵的特征值}
    设 可逆阵$\bm{A}$的特征值为$\lambda_1,\lambda_2,\cdots,\lambda_n$,则$\bm{A}^{-1}$的特征值为$\lambda_1^{-1},\lambda_2^{-1},\cdots,\lambda_n^{-1}$.
\end{proposition}
\begin{Proof}
因为存在非异阵$\bm{P}\in M_n\left(
    \mathbb{K}
\right)$使得
\[
\bm{P}^{-1}\bm{AP}=\begin{pmatrix}
    \lambda_1&\cdots&*\\
    \vdots&\ddots&\vdots\\
    \bm{O}&\cdots&\lambda_n
\end{pmatrix}
\]因为非异上三角阵的逆阵仍为非异上三角阵,且主对角线元素的逆等于原主对角线元素的逆,所以
\[
\bm{P}^{-1}\bm{A}^{-1}\bm{P}=\begin{pmatrix}
    \lambda_1^{-1}&\cdots&*\\
    \vdots&\ddots&\vdots\\
    \bm{O}&\cdots&\lambda_n^{-1}
\end{pmatrix}\qedhere
\]
\end{Proof}
\end{formal}
\begin{formal}
\begin{proposition}[伴随阵的特征值]\label{prop:伴随阵的特征值}
    设$\bm{A}\in M_n\left(
        \mathbb{K}
    \right)$的特征值为$\lambda_1,\lambda_2,\cdots,\lambda_n$,则$\bm{A}^*$的特征值为
    $\displaystyle
    \prod_{i\neq 1}\lambda_i,\prod_{i\neq 2}\lambda_i,\cdots,\prod_{i\neq n}\lambda_i
    $.
\end{proposition}
\begin{Proof}
注意到上三角阵的伴随阵仍为上三角阵,且主对角线元素的伴随等于其余元素的乘积即可.
\end{Proof}
\end{formal}
\newpage
\section{对角化}
\subsection{前提}
\begin{formal}
\begin{definition}[几何意义的对角化]\label{def:几何意义的对角化}
    设$\bm{\varphi}\in \mathcal{L}\left(
        V_{\mathbb{K}}^n
    \right)$,若$\bm{\varphi}$在某组基下的表示矩阵为对角阵,则称$\bm{\varphi}$可对角化.
\end{definition}
\end{formal}
\begin{formal}
\begin{definition}[代数意义的对角化]\label{def:代数意义的对角化}
    设$\bm{A}\in M_n\left(
        \mathbb{K}
    \right)$,若$\bm{A}\approx \bm{\varLambda}$则称$\bm{A}$可对角化.
\end{definition}
\end{formal}
\begin{red}
\begin{remark}
    前文已述,一般来说是在复数域上进行考虑的,有时为了便于区分,也称为复可对角化.
\end{remark}
\end{red}
\begin{formal}
\begin{theorem}[几何意义的可对角化条件]\label{thm:几何意义的可对角化条件}
    设$V$是$\mathbb{K}$上的$n$维线性空间,$\bm{\varphi}\in\mathcal{L}\left(V\right)$.则$\bm{\varphi}$可对角化当且仅当$\bm{\varphi}$有$n$个线性无关的特征向量.
\end{theorem}
\begin{Proof}
先考虑必要性,设$\bm{\varphi}$可对角化,即存在$V$的一组基$\left\{
    \bm{e}_1,\bm{e}_2,\cdots,\bm{e}_n
\right\}$使得$\bm{\varphi}$在这组基下的表示矩阵为$\mathrm{diag}\,\left\{
    \lambda_1,\lambda_2,\cdots,\lambda_n
\right\}$,则有
\[
\left(
    \bm{\varphi}\left(\bm{e}_1\right),
    \bm{\varphi}\left(\bm{e}_2\right),
    \cdots,
    \bm{\varphi}\left(\bm{e}_n\right)
\right)=
\left(
    \bm{e}_1, \bm{e}_2,\cdots,\bm{e}_n
\right)\begin{bmatrix}
    \lambda_1&&&\\
    &\lambda_2&&\\
    &&\ddots&\\
    &&&\lambda_n
\end{bmatrix}
\]于是$\bm{\varphi}\left(\bm{e}_i\right)=\lambda_i\bm{e}_i,\forall 1\leqslant i\leqslant n.$于是$\lambda_1,\lambda_2,\cdots,\lambda_n$正是
$\bm{\varphi}$的特征值,而$\bm{e}_1,\bm{e}_2,\cdots,\bm{e}_n$正是对应于这$n$个特征值的特征向量,并且它们作为基向量一定是线性无关的.必要性证毕.

再考虑充分性,设$\bm{\varphi}$有$n$个线性无关的特征向量$\bm{e}_1,\bm{e}_2,\cdots,\bm{e}_n$,则它们构成$V$的一组基,且$\bm{\varphi}\left(\bm{e}_i\right)=\lambda_i\bm{e}_i,\forall 1\leqslant i\leqslant n$.于是
\[
\left(
    \bm{\varphi}\left(\bm{e}_1\right),
    \bm{\varphi}\left(\bm{e}_2\right),
    \cdots,
    \bm{\varphi}\left(\bm{e}_n\right)
\right)=
\left(
    \bm{e}_1, \bm{e}_2,\cdots,\bm{e}_n
\right)\begin{bmatrix}
    \lambda_1&&&\\
    &\lambda_2&&\\
    &&\ddots&\\
    &&&\lambda_n
\end{bmatrix}
\]
即$\bm{\varphi}$在这组基下的表示矩阵为对角阵,故$\bm{\varphi}$可对角化.

证毕.\qedhere
\end{Proof}
\end{formal}
\begin{formal}
\begin{theorem}[代数意义的可对角化条件]\label{thm:代数意义的可对角化条件}
    设$\bm{A}\in M_n\left(
        \mathbb{K}
    \right)$则$\bm{A}$可复对角化当且仅当
    $\bm{A}$有$n$个线性无关的特征向量.
\end{theorem}
\end{formal}
\begin{red}
\begin{remark}
    \cref{thm:代数意义的可对角化条件}显然不是任一矩阵均满足的.
\end{remark}
\end{red}
\begin{brown}
\begin{example}
    考虑矩阵
    \[
    \begin{bmatrix}
        1&1\\
        0&1
    \end{bmatrix}
    \]其特征值为$1,1$,但只有一个线性无关的特征向量$k\begin{pmatrix}
        1\\0
    \end{pmatrix}\left(
        k\in\mathbb{K}\backslash \left\{0\right\}
    \right)$,故不可对角化.

    或者考虑反证法,若$\bm{A}$可对角化,则存在非异阵$\bm{P}\in M_2\left(
        \mathbb{K}
    \right)$使得
    $\bm{P}^{-1}\bm{AP}=\bm{I}_2$于是
    $\bm{A}=\bm{P}\bm{I}_2\bm{P}^{-1}=\bm{I}_2$,矛盾.
\end{example}
\end{brown}
\begin{formal}
\begin{theorem}[特征子空间的和]\label{thm:特征子空间的和}
    设$\bm{\varphi}\in
    \mathcal{L}\left(
        V_{\mathbb{K}}^n
            \right)
    $,$\lambda_1,\lambda_2,\cdots,\lambda_k$是$\bm{\varphi}$的不同特征值,$V_1,V_2,\cdots,V_k$是对应的特征子空间,则
    \[
     V_1+V_2+\cdots+V_k=V_1\oplus 
        V_2\oplus\cdots\oplus V_k
    \]
\end{theorem}
\begin{Proof}
对$k$进行归纳,$k=1$显然.设不同特征值的个数$<k$时成立,下面考虑其等于$k$时的情况.此时有
\[
V_1+V_2+\cdots+V_{k-1}=V_1\oplus 
        V_2\oplus\cdots\oplus V_{k-1}
        \]即$V_i\cap
        \left(
            V_1+V_2+\cdots+V_{i-1}
        \right)=0,\forall 2\leqslant i\leqslant k-1$.要证
        \[
        V_1+V_2+\cdots+V_k=V_1\oplus
        V_2\oplus\cdots\oplus V_k
        \]只需证$V_k\cap
        \left(
            V_1+V_2+\cdots+V_{k-1}
        \right)=0$.

        任取$\bm{v}\in V_k\cap\left(
            V_1+V_2+\cdots+V_{k-1}
        \right)$则
        \[
        V_k \ni \bm{v} = 
        \bm{v}_1+\bm{v}_2+\cdots+\bm{v}_{k-1}
        \]其中$\bm{v}_i\in V_i,\forall 1\leqslant i\leqslant k -1.$两边同时作用$\bm{\varphi}$,则
        \[
         \lambda_k\bm{v} = 
        \lambda_1\bm{v}_1+\lambda_2\bm{v}_2+\cdots+\lambda_{k-1}\bm{v}_{k-1}
        \]又因为
        \[
        \lambda_k\bm{v} = 
        \lambda_k\bm{v}_1+\lambda_k\bm{v}_2+\cdots+\lambda_{k}\bm{v}_{k-1}
        \]即有
        \[
         \bm{0} = 
        \left(
            \lambda_1-\lambda_k
        \right)\bm{v}_1+\left(
            \lambda_2-\lambda_k
        \right)\bm{v}_2+\cdots+\left(
            \lambda_{k-1}-\lambda_k
        \right)\bm{v}_{k-1}
        \]
        由直和假设知零向量分块表示唯一,又因为$\lambda_i\neq \lambda_j,\forall i\neq j$,故$\bm{v}_1=\bm{v}_2=\cdots=\bm{v}_{k-1}=\bm{0}$,即$\bm{v}=\bm{0}$,从而
        \[
        V_1+V_2+\cdots+V_k=V_1\oplus
        V_2\oplus\cdots\oplus V_k
        \]
        得证.\qedhere
\end{Proof}
\end{formal}
\begin{green}
\begin{corollary}[不同特征值的特征向量]\label{cor:不同特征值的特征向量}
    不同特征值的特征向量必定线性无关.
\end{corollary}
\begin{Proof}
设$\lambda_1,\lambda_2,\cdots,\lambda_k$是不同的特征值,$\bm{v}_1,\bm{v}_2,\cdots,\bm{v}_k$是对应的特征向量.设
\[
c_1\bm{v}_1+c_2\bm{v}_2+\cdots+c_k\bm{v}_k=\bm{0}
\]
由上定理及直和的零向量分块表示唯一且$\bm{v}_i\neq \bm{0}$知$c_1=c_2=\cdots=c_k$.
\end{Proof}
\end{green}
\begin{formal}
\begin{theorem}[特征值数与可对角化]\label{thm:特征值数与可对角化}
    设$\bm{\varphi}\in \mathcal{L}\left(
        V_{\mathbb{K}}^n
    \right)$有$n$个不同的特征值,则
    $\bm{\varphi}$可对角化.
\end{theorem}
\begin{Proof}
设$\bm{\varphi}$的特征值为$\lambda_1,\lambda_2,\cdots,\lambda_n$,由于它们不同,故对应的特征向量$\bm{v}_1,\bm{v}_2,\cdots,\bm{v}_n$线性无关,于是$\bm{\varphi}$有$n$个线性无关的特征向量,故可对角化.
\end{Proof}
\end{formal}
\begin{red}
\begin{remark}
    \cref{thm:特征值数与可对角化}是一个充分条件,但不是必要条件.
\end{remark}
\end{red}
\begin{brown}
\begin{example}
    设$\bm{\varphi}=c\bm{I}_V,\bm{A}=c\bm{I}_n$是可对角化的但特征值全为$c$.
\end{example}
\end{brown}
\begin{red}
\begin{remark}
    实际上,任给一个矩阵,其可对角化的概率要远大于不可对角化的概率.

    并且可以用可对角化的矩阵取去逼近不可对角化的矩阵.
\end{remark}
\end{red}
\begin{formal}
\begin{theorem}[可对角化的条件]\label{thm:可对角化的条件}
    设$V$是$\mathbb{C}$上的$n$维线性空间,$
     \bm{\varphi}\in \mathcal{L}\left(
        V
    \right)
    $,设全体不同特征值$\lambda_1,\lambda_2,\cdots,\lambda_k$,$V_1,V_2,\cdots,V_k$为
    对应的特征子空间,则$\bm{\varphi}$可对角化当且仅当
    \[
    V=V_1\oplus V_2\oplus\cdots\oplus V_k
    \]
\end{theorem}
\begin{Proof}
考虑充分性,设
\[
V=V_1\oplus V_2\oplus\cdots\oplus V_k
\]取$V_1$的基$\left\{
    \bm{v}_{11},\bm{v}_{12},\cdots,\bm{v}_{1n_1}
\right\}$,$V_2$的基$\left\{
    \bm{v}_{21},\bm{v}_{22},\cdots,\bm{v}_{2n_2}
\right\}$,$\cdots$,$V_k$的基$\left\{
    \bm{v}_{k1},\bm{v}_{k2},\cdots,\bm{v}_{kn_k}
\right\}$,则$
\left\{
    \bm{v}_{11},\bm{v}_{12},\cdots,\bm{v}_{1n_1},\bm{v}_{21},\bm{v}_{22},\cdots,\bm{v}_{2n_2},\cdots,\bm{v}_{k1},\bm{v}_{k2},\cdots,\bm{v}_{kn_k}
\right\}
$是$V$的一组基.于是$\bm{\varphi}$有$n$个线性无关的特征向量,故可对角化.

再考虑必要性,设$\bm{\varphi}$有$n$个线性无关的特征向量$\bm{e}_1,\bm{e}_2,\cdots,\bm{e}_n$,这也是$V$的一组基,于是$\forall \bm{\alpha}\in V$
\[
\bm{\alpha}=a_1\bm{e}_1 + a_2\bm{e}_2+\cdots+a_n\bm{e}_n
\]
进行适当调序可以保证$a_{i_1}\bm{e}_{i_1},a_{i_1}\bm{e}_{i_r}\cdots,a_{i_r}\bm{e}_{i_r}\in V_r\left(
    1   \leqslant r\leqslant k
\right)$,于是
\[
V=V_1+ V_2+\cdots+V_k
\]又考虑到不同特征值的特征子空间的和一定是直和,故
\[
V=V_1\oplus V_2\oplus\cdots\oplus V_k
\qedhere
\]
\end{Proof}
\end{formal}
\subsection{重数}
\begin{formal}
\begin{definition}[几何重数与代数重数]\label{def:几何重数与代数重数}
    设$\lambda_0$是$\bm{\varphi}$的特征值,$V_0$是其特征子空间,则称$\dim\,V_0$为$\lambda_0$的度数或几何重数.

    $\lambda_0$作为$\bm{\varphi}$的特征多项式的根的重数称为$\lambda_0$的重数或代数重数.
\end{definition}
\end{formal}
\begin{red}
\begin{remark}
    几何重数与代数重数不见得相等.由线性映射维数公式知,几何重数
    \[
    \dim\,V_0=\dim\,\mathrm{Ker}\,\left(
        \lambda_0\bm{I}_V-\bm{\varphi}
    \right)=n-\mathrm{r}\,\left(
        \lambda_0\bm{I}_V-\bm{\varphi}
    \right)
    \]
    代数重数取决于特征多项式$\left|
    \lambda\bm{I}_V-\bm{\varphi}
    \right|$的因式分解.
\end{remark}
\end{red}
\begin{green}
\begin{lemma}[代数重数与几何重数的关系]\label{lem:代数重数与几何重数的关系}
    设$
    \bm{\varphi}\in \mathcal{L}\left(
        V_{\mathbb{K}}^n
    \right)
    $,$\lambda_0$是$\bm{\varphi}$的一个特征值,则$\lambda_0$的几何重数一定小于等于其代数重数.
\end{lemma}
\begin{Proof}
设几何重数$t=\dim\,V_0$,代数重数为$m$.

取$V_0$的一组基$\left\{
    \bm{e}_1    ,\bm{e}_2,\cdots,\bm{e}_t
\right\}$扩张为$V$的一组基$\left\{
    \bm{e}_1,\bm{e}_2,\cdots,\bm{e}_t,\bm{e}_{t+1},\cdots,\bm{e}_n  
\right\}$,则有
\begin{align*}
   & \left(
     \bm{\varphi}\left(
            \bm{e}_1 
     \right),
        \bm{\varphi}\left(
                \bm{e}_2
        \right),
        \cdots,
        \bm{\varphi}\left(
                \bm{e}_t
        \right),
        \bm{\varphi}\left(
                \bm{e}_{t+1}
        \right),
        \cdots,
        \bm{\varphi}\left(
                \bm{e}_n
                \right)
\right)\\
=&\left(
    \bm{e}_1,\bm{e}_2,\cdots,\bm{e}_t,\bm{e}_{t+1},\cdots,\bm{e}_n
\right)\begin{bmatrix}
    \lambda_0\bm{I}_t&\bm{C}\\
    \bm{O}&\bm{B}
\end{bmatrix}
\end{align*}
不妨设$\bm{A}=\begin{bmatrix}
    \lambda_0\bm{I}_t&\bm{C}\\
    \bm{O}&\bm{B}
\end{bmatrix}$,于是特征多项式
\begin{align*}
    \left|\lambda\bm{I}_V-\bm{\varphi}\right|&=
    \left|
       \lambda \bm{I}_n-\bm{A}
    \right|\\
    &=
    \begin{vmatrix}
        \left(
            \lambda-\lambda_0
        \right)\bm{I}_t&-\bm{C}\\
        \bm{O}&\lambda\bm{I}_{n-t}-\bm{B}
    \end{vmatrix}
\end{align*}
由$\textup{Laplace}$定理知上式等于$\left(
    \lambda-\lambda_0
\right)^{\prime}\left|
    \lambda\bm{I}_{n-t}-\bm{B}
\right|$.于是$m\geqslant t.$
\end{Proof}
\end{green}
\begin{formal}
\begin{definition}[完全特征向量系]\label{def:完全特征向量系}
    设
    $\bm{\varphi}\in \mathcal{L}\left(
        V_{\mathbb{K}}^n
    \right)
    $,若对$\bm{\varphi}$的任一特征值$\lambda_0$,其几何重数等于代数重数,则称$\bm{\varphi}$有完全的特征向量系.
\end{definition}
\end{formal}
\begin{formal}
\begin{theorem}[可对角化的条件]\label{thm:可对角化的条件2}
    设
    $\bm{\varphi}\in \mathcal{L}\left(
        V_{\mathbb{K}}^n
    \right)
    $,则$\bm{\varphi}$可对角化当且仅当$\bm{\varphi}$有完全的特征向量系.
\end{theorem}
\begin{Proof}
只要证明$\bm{\varphi}$有完全的特征向量系当且仅当
\[
V= V_1\oplus V_2\oplus\cdots\oplus V_k
\]
其中,$\lambda_1,\lambda_2,\cdots,\lambda_k$是$\bm{\varphi}$的全体不同特征值,$V_1,V_2,\cdots,V_k$是对应的特征子空间.

先考虑充分性,设$\lambda_i$的几何重数
$t_i=\dim\,V_i\left(
    \forall 1\leqslant i\leqslant k 
\right)$和代数重数$m_i\left(
    \forall 1\leqslant i\leqslant k
\right)$.于是
$t_i\leqslant m_i\left(
    \forall 1\leqslant i\leqslant k
\right)$且$m_1+m_2+\cdots+m_k=n.$因为
\begin{align*}
    n&=\dim\,V=\dim\,\left(
        V_1\oplus   V_2\oplus\cdots\oplus V_k
    \right)\\
    &=\dim\,V_1+\dim\,V_2+\cdots+\dim\,V_k\\
    &=t_1+t_2+\cdots+t_k\leqslant m_1+m_2+\cdots+m_k=n\\
    &\Longrightarrow 
    m_i=t_i\left(
        \forall 1\leqslant i\leqslant k
    \right)
\end{align*}
于是$\bm{\varphi}$有完全的特征向量系.

在考虑必要性,即$\bm{\varphi}$有完全的特征向量系即$t_i=m_i\left(
    \forall 1\leqslant i\leqslant k
\right)$,则
\begin{align*}
    \dim\,\left(
       V_1\oplus V_2\oplus\cdots\oplus V_k   
\right)&=\dim\,V_1+\dim\,V_2+\cdots+\dim\,V_k
\\
&=t_1+t_2+\cdots+t_k\\
&=m_1+m_2+\cdots+m_k\\
&=n=\dim\,V\\
&\Longrightarrow V=V_1\oplus V_2\oplus\cdots\oplus V_k\qedhere
\end{align*}
\end{Proof}
\end{formal}
\begin{brown}
\begin{example}
考虑矩阵
\newpage
\[
\bm{A}=\begin{bmatrix}
    1&1\\
    0&1
\end{bmatrix}
\]对于其特征值$1$,代数重数为$2$,几何重数为$1$,故没有完全的特征向量系,因此不可对角化.
\end{example}
\end{brown}
\begin{brown}
\begin{example}
    若已知$\bm{A}$可对角化,如何求得$\bm{P}$使得$\bm{P}^{-1}\bm{AP}$为对角阵.
\end{example}
\begin{solution}
    令$\bm{P}=\left(
        \bm{\alpha}_1,\bm{\alpha}_2,\cdots,\bm{\alpha}_n
    \right)$,其为非异阵即$\bm{\alpha}_1,
    \bm{\alpha}_2,\cdots,\bm{\alpha}_n$线性无关.若
    \[
    \bm{P}^{-1}\bm{AP}=\begin{pmatrix}
        \lambda_1&&&\\
        &\lambda_2&&\\
        &&\ddots&\\
        &&&\lambda_n
    \end{pmatrix}
    \]则
    \[
    \bm{AP}=
    \bm{P}\begin{pmatrix}
        \lambda_1&&&\\
        &\lambda_2&&\\
        &&\ddots&\\
        &&&\lambda_n
    \end{pmatrix}
    \]即$\bm{A\alpha}_1=\lambda_1\bm{\alpha}_1,
    \bm{A\alpha}_2=\lambda_2\bm{\alpha}_2,\cdots,
    \bm{A\alpha}_n=\lambda_n\bm{\alpha}_n
    $.即$\bm{P}$的$n$个列向量就是$\bm{A}$的$n$个线性无关的特征向量.
    
    并且由于特征向量的线性组合仍为特征向量,故$\bm{P}$的选取一定不是唯一的.
\end{solution}
\end{brown}
\begin{red}
\begin{remark}
    $\bm{P}$的第$i$个列向量对应于$\bm{A}$的第$i$个特征值.
\end{remark}
\end{red}
\begin{brown}
\begin{example}
    考虑矩阵\newpage
    \[
    \bm{A}=\begin{bmatrix}
        1&0&0\\-2&5&-2\\-2&4&-1
    \end{bmatrix}
    \]
\end{example}
\begin{Proof}
先计算各代数重数与几何重数
\begin{align*}
    \left|
    \lambda\bm{I}_3-\bm{A}
    \right|&=\begin{vmatrix}
        \lambda-1&0&0\\2&\lambda-5&2\\2&-4&\lambda+1
    \end{vmatrix}\\
    &=\left(
        \lambda-1
    \right)^2\left(
        \lambda-3
    \right)
\end{align*}于是特征值$\lambda_1=1$,代数重数为$2$,解方程组得
\begin{align*}
    \begin{bmatrix}
        0&0&0\\2&-4&2\\2&-4&2
    \end{bmatrix}\longrightarrow
    \begin{bmatrix}
        1&-2&1\\0&0&0\\0&0&0
    \end{bmatrix}
\end{align*}基础解系为$\begin{pmatrix}
    2\\1\\0
\end{pmatrix},\begin{pmatrix}
    -1\\0\\1
\end{pmatrix}$,故几何重数为$2$.

特征值$\lambda_2=3$,代数重数为$1$,解方程组得
\begin{align*}
    \begin{bmatrix}
        2&0&0\\2&-2&2\\2&-4&4
    \end{bmatrix}\longrightarrow
    \begin{bmatrix}
        1&0&0\\0&1&-1\\0&0&0
    \end{bmatrix}
\end{align*}基础解系为$\begin{pmatrix}
    0\\1\\1
\end{pmatrix}$,故几何重数为$1$（事实上,可以证明,代数重数为$1$时几何重数必定为$1$）.

于是$\bm{A}$一定可对角化且
\[
\bm{P}=\begin{bmatrix}
    2&-1&0\\1&0&1\\0&1&1
\end{bmatrix}
,
\bm{P}^{-1}\bm{AP}=\begin{bmatrix}
    1&0&0\\0&1&0\\0&0&3
\end{bmatrix}
\]
\end{Proof}
\end{brown}
\begin{brown}
\begin{example}
    求可对角化矩阵$\bm{A}$的幂$\bm{A}^m$.
\end{example}
\begin{solution}
    显然
    \[
    \bm{A}^m=
    \bm{P}\begin{bmatrix}
        \lambda_1^m&&&\\
        &\lambda_2^m&&\\
        &&\ddots&\\
        &&&\lambda_n^m
    \end{bmatrix}\bm{P}^{-1}
    \]
\end{solution}
\end{brown}
\begin{brown}
\begin{example}
    设$\bm{A}=\begin{bmatrix}
        1&0\\1&-2
    \end{bmatrix}$,求$\bm{A}^{m}.$
\end{example}
\begin{solution}
    注意到这是一个下三角阵,特征值为$1$和$-2$,于是一定有两个线性无关的特征向量,故一定可对角化.求出\[
    \bm{P}=\begin{bmatrix}
        3&0 \\ 1&1
    \end{bmatrix}
    \]使得
    \[
    \bm{P}^{-1}\bm{AP}=\begin{bmatrix}
        1&0\\0&-2
    \end{bmatrix}
    \]于是
    \begin{align*}
    \bm{A}^m&=\bm{P}\begin{bmatrix}
        1^m&0\\0&(-2)^m
    \end{bmatrix}\bm{P}^{-1}\\
    &=\begin{bmatrix}
        1&0\\
        \frac{1}{3}\left(
            1   -(-2)^m
        \right)&(-2)^m
    \end{bmatrix}
\end{align*}
\end{solution}
\end{brown}
\begin{red}
\begin{remark}
    事实上,考虑几何重数还可以有另外的方法,这在前文已述.因为
\[
V_{\lambda_0}=\left\{
  \bm{v}\in \mathbb{C}^n\mid
  \bm{A}\bm{v}=\lambda_0\bm{v}   
\right\}=\left(
    \lambda_0\bm{I}_n-\bm{A}
\right)\bm{x}=\bm{0}\textup{的解空间}
\]故
几何重数$\dim\,V_{\lambda_0}=n-\mathrm{r}\left(
    \lambda_0\bm{I}_n-\bm{A}
\right)$.
\end{remark}
\end{red}
\newpage
\section{极小多项式与Cayley-Hamilton定理}
\subsection{极小多项式}
    设$\bm{A}\in M_n\left(
        \mathbb{K}
    \right)$,因为
    \[
    \dim_{\mathbb{K}}\,M_n\left(
        \mathbb{K}
    \right)=n^2
    \]故$\bm{A}^{n^2},\bm{A}^{n^2-1},\cdots,\bm{A},\bm{I}_n$这$n^2+1$个矩阵必定线性相关.
    
    即存在不全为$0$的数$c_{n^2},c_{n^2-1},\cdots,c_1,c_0\in\mathbb{K}\,\textup{s.t.}$
    \[
    c_{n^2}\bm{A}^{n^2}+c_{n^2-1}\bm{A}^{n^2-1}+\cdots+c_1\bm{A}+c_0\bm{I}_n=\bm{O}
    \]即$\bm{A}$适合多项式$g\left(x\right)=
    c_{n^2}x^{n^2}+c_{n^2-1}x^{n^2-1}+\cdots+c_1x+c_0\neq 0\in\mathbb{K}\left[x\right].$

    定义$S=\left\{
        f\left(x\right)\in\mathbb{K}\left[x\right]\mid 
        f\left(x\right)\neq 0 \textup{且} f\left(\bm{A}\right)=\bm{O}
    \right\}\neq \varnothing $,则$\displaystyle
    k=\min_{f\left(x\right)\in S}\deg\,f\left(x\right)$必定存在,即$\exists h\left(x\right)\in S\,\textup{s.t.}\deg\,h\left(x\right)=k.$首一化得到$m\left(x\right)$,即有首一多项式$0\neq m\left(x\right)\in \mathbb{K}\left[x\right]\,\textup{s.t.}m\left(\bm{A}\right)=\bm{O}$且有次数最小性.
\begin{formal}
\begin{definition}[极小多项式]\label{def:极小多项式}
    设$\bm{A}\in M_n\left(
        \mathbb{K}
    \right),m\left(x\right)\in \mathbb{K}\left[x\right]$为首一多项式且$m\left(\bm{A}\right)=\bm{O}$.若$m\left(x\right)$是$\bm{A}$适合的所有非零多项式中次数最小者,则称为$\bm{A}$的一个极小多项式（最小多项式）.
\end{definition}
\end{formal}
\begin{red}
\begin{remark}
    极小多项式一定存在.
\end{remark}
\end{red}
\begin{green}
\begin{lemma}[整除]\label{lem:整除}
    设$\bm{A}$适合$f\left(x\right)\neq 0$,$m\left(x\right)$是$\bm{A}$的极小多项式则$m\left(x\right)\mid f\left(x\right).$
\end{lemma}
\begin{Proof}
考虑带余除法$f\left(x\right)=m\left(x\right)q\left(x\right)+r\left(x\right),\deg\,r\left(x\right)<\deg\,m\left(x\right)$,考虑反证法,即设$r\left(x\right)\neq 0$,上式代入$\bm{A}$,于是$r\left(\bm{A}\right)=\bm{O}$,这与$m\left(x\right)$的次数最小性矛盾.
\end{Proof}
\end{green}
\begin{formal}
\begin{proposition}[极小多项式的唯一性]\label{prop:极小多项式的唯一性}
    设$\bm{A}\in M_n\left(\mathbb{K}\right)$,其极小多项式存在且唯一.
\end{proposition}
\begin{Proof}
设$m\left(x\right),g\left(x\right)$均为$\bm{A}$的极小多项式,只需证$m\left(x\right)=g\left(x\right)$.事实上,$m\left(x\right)\mid g\left(x\right),g\left(x\right)\mid m\left(x\right)$即$\exists 0\neq c\in\mathbb{K}\,\textup{s.t.}g\left(x\right)=cm\left(x\right)$但极小多项式均首一于是$c=1$.
\end{Proof}
\end{formal}
\begin{brown}
\begin{example}
    纯量矩阵$\bm{A}=c\bm{I}_n$的极小多项式为$m\left(x\right)=x-c$.
\end{example}
\end{brown}
\begin{brown}
\begin{example}
    考虑矩阵$\bm{A}=\begin{bmatrix}
        0&1\\0&0
    \end{bmatrix}$其满足$\bm{A}^2=\bm{O}$即$\bm{A}$适合多项式$f\left(x\right)=x^2$,故其极小多项式为必定整除$f\left(x\right)=x^2$即$m\left(x\right)=x$或$m\left(x\right)=x^2$于是$m\left(x\right)=x^2.$
\end{example}
\end{brown}
\begin{formal}
\begin{proposition}[相似矩阵的极小多项式]\label{prop:相似矩阵的极小多项式}
    相似矩阵具有相同的极小多项式.
\end{proposition}
\begin{Proof}
设$\bm{B}=\bm{P}^{-1}\bm{AP}$,$\bm{A}$的极小多项式为$m\left(x\right)$,$\bm{B}$的极小多项式为$g\left(x\right)$,则考虑到相似关系与多项式的相容有
\[
g\left(\bm{A}\right)=g\left(\bm{P}\bm{BP}^{-1}\right)=\bm{P}g\left(\bm{B}\right)\bm{P}^{-1}=\bm{O}
\]于是$m\left(x\right)\mid g\left(x\right)$,同理得证.
\end{Proof}
\end{formal}
\begin{formal}
\begin{proposition}[分块对角矩阵的极小多项式]\label{prop:分块对角矩阵的极小多项式}
    设分块对角阵
    \[
    \bm{A}=\begin{bmatrix}
        \bm{A}_1&&&\\
        &\bm{A}_2&&\\
        &&\ddots&\\
        &&&\bm{A}_k  
    \end{bmatrix}
    \]
    其中$\bm{A}_i\left(
        1\leqslant i\leqslant k
    \right)$为方阵,则$\bm{A}$的极小多项式为诸$\bm{A}_i$的极小多项式的最小公倍式即
   $\left[
    m_1\left(x\right),m_2\left(x\right),\cdots,m_k\left(x\right)
   \right]$其中$m_i\left(x\right)$为$\bm{A}_i$的极小多项式.
\end{proposition}
\begin{Proof}
令$g\left(x\right)=\left[
    m_1\left(x\right),m_2\left(x\right),\cdots,m_k\left(x\right)
\right]$只需证$m\left(x\right)=g\left(x\right)$.显然$m_i\left(x\right)\mid g\left(x\right)$,于是$g\left(\bm{A}_i\right)=\bm{O}$.因为\[
g\left(\bm{A}\right)=\begin{bmatrix}
    g\left(\bm{A}_1\right)&&&\\
    &g\left(\bm{A}_2\right)&&\\
    &&\ddots&\\
    &&&g\left(\bm{A}_k\right)
\end{bmatrix}=\bm{O}
\]于是$m\left(x\right)\mid g\left(x\right)$.

另一方面
\[
\bm{O}=m\left(\bm{A}\right)=\begin{bmatrix}
    m\left(\bm{A}_1\right)&&&\\
    &m\left(\bm{A}_2\right)&&\\
    &&\ddots&\\
    &&&m\left(\bm{A}_k\right)
\end{bmatrix}
\]于是$m\left(\bm{A}_i\right)=\bm{O}$故$m_i\left(x\right)\mid m\left(x\right),\forall 
1\leqslant i\leqslant k$.由于$g\left(x\right)$是最小公倍式则$g\left(x\right)\mid m\left(x\right)$.类似得证.
\end{Proof}
\end{formal}
\begin{brown}
\begin{example}\label{ex:极小多项式无重根等价于可对角化}
    设$\bm{A}$的全体不同特征值$\lambda_1,
    \lambda_2,\cdots,\lambda_k$,证明：若$\bm{A}$可对角化,则其极小多项式为
    $m\left(x\right)=\left(
        x-\lambda_1
    \right)\left(
        x-\lambda_2
    \right)\cdots\left(
        x-\lambda_k
    \right)$.
\end{example}
\begin{Proof}
因为$\exists \textup{非异阵}\bm{P}\,\textup{s.t.}$
\[
\bm{P}^{-1}\bm{AP}=\begin{pmatrix}
    \lambda_1\bm{I}&&&\\
    &\lambda_2\bm{I}&&\\
    &&\ddots&\\
    &&&\lambda_k\bm{I}
\end{pmatrix}=\bm{B}
\]于是$\bm{A}$的极小多项式也为$\bm{B}$的极小多项式,而由上命题知后者等于\[\left[
    x-\lambda_1,x-\lambda_2,\cdots,x-\lambda_k
\right]=\left(
    x-\lambda_1
\right)
    \left(
        x-\lambda_2
    \right)\cdots\left(
        x-\lambda_k
    \right)\qedhere
\]
\end{Proof}
\end{brown}
\begin{red}
\begin{remark}
    于是可对角化的矩阵的特征值一定是极小多项式的根.

    事实上,对于任一矩阵也是如此.
\end{remark}
\end{red}
\begin{red}
    \begin{remark}
        后续会证明：若矩阵的极小多项式无重根,则矩阵一定可对角化.

        即一个矩阵可对角化当且仅当其极小多项式无重根.
    \end{remark}
\end{red}
\begin{green}
\begin{lemma}[特征值与根]\label{lem:特征值与根}
    设$\bm{A}\in M_n\left(
        \mathbb{K}
    \right)$的极小多项式为$m\left(x\right)$,$\lambda_0$是$\bm{A}$的一个特征值,则
    \[
    \left(
        x-\lambda_0
    \right)\mid m\left(x\right)
    \]
\end{lemma}
\begin{Proof}
因为$m\left(\bm{A}\right)=\bm{O}$,因为一个矩阵适合一个多项式则其特征值均适合该多项式即$m\left(\lambda_0\right)=0$.由\cref{thm:余数定理}余数定理立得.
\end{Proof}
\end{green}
\subsection{Cayley-Hamilton定理}
\begin{formal}
\begin{proposition}[简单版]\label{prop:简单版}
    设
    \[
    \bm{A}=\begin{pmatrix}
        \lambda_1&a_{12}&\cdots&a_{1n}\\
        &\lambda_2&\cdots&a_{2n}\\
        &&\ddots&\vdots\\
        &&&\lambda_n
    \end{pmatrix}
    \]则$\left(
        \bm{A}-\lambda_1\bm{I}_n
    \right)
    \left(
        \bm{A}-\lambda_2\bm{I}_n
    \right)\cdots\left(
        \bm{A}-\lambda_n\bm{I}_n
    \right)=\bm{O}.
    $即上三角阵$\bm{A}$适合自己的特征多项式$f\left(\lambda\right)=\left(\lambda-\lambda_1\right)
    \left(\lambda-\lambda_2\right)\cdots\left(\lambda-\lambda_n\right).
    $
\end{proposition}
\begin{Proof}
取标准单位列向量,于是
\[
\bm{Ae}_1=\lambda_1\bm{e}_1,\bm{Ae}_2=a_{12}\bm{e}_1+\lambda_2\bm{e}_2,\cdots,\bm{Ae}_n=a_{1n}\bm{e}_1+a_{2n}\bm{e}_2+\cdots+\lambda_n\bm{e}_n
\]要证明
\[
\left(
    \bm{A}-\lambda_1\bm{I}_n   
\right)
\left(
    \bm{A}-\lambda_2\bm{I}_n
\right)\cdots\left(
    \bm{A}-\lambda_n\bm{I}_n
\right)=\bm{O}
\]只要证
\[
\left(
    \bm{A}-\lambda_1\bm{I}_n
\right)\left(
    \bm{A}-\lambda_2\bm{I}_n
\right)
\cdots\left(
    \bm{A}-\lambda_n\bm{I}_n
\right)\bm{e}_i=\bm{0},\forall 1\leqslant i\leqslant n
\]甚至,只要证明
\[
\left(
    \bm{A}-\lambda_1\bm{I}_n
\right)\left(
    \bm{A}-\lambda_2\bm{I}_n
\right)\cdots\left(
    \bm{A}-\lambda_i\bm{I}_i
\right)\bm{e}_i=\bm{0},\forall 1\leqslant i\leqslant n
\]即可（这是因为我们有$f\left(\bm{A}\right)g\left(\bm{A}\right)=g\left(\bm{A}\right)f\left(\bm{A}\right)$）.

于是考虑归纳法证明,对$i$进行归纳,一方面$i=1$时,显然$\left(
    \bm{A}-\lambda_1\bm{I}_n
\right)\bm{e}_1=\bm{0}$.设$<i$时结论成立,则证$\left(
    \bm{A}-\lambda_1\bm{I}_n
\right)
\left(
    \bm{A}-\lambda_2\bm{I}_n
\right)\cdots\left(
    \bm{A}-\lambda_i\bm{I}_n
\right)\bm{e}_i=\bm{0}
$
\begin{align*}
    LHS &= \left(
        \bm{A}-\lambda_1\bm{I}_n
    \right)
    \left(
        \bm{A}-\lambda_2\bm{I}_n
    \right)\cdots\left(
        \bm{A}-\lambda_{i-1}\bm{I}_n
    \right)\left(
        a_{1i}\bm{e}_1+a_{2i}\bm{e}_2+\cdots+a_{i-1,i}\bm{e}_{i-1}
    \right)
\end{align*}
展开后由归纳假设知上式等于$\bm{0}$,故得证.
\end{Proof}
\end{formal}
\begin{formal}
\begin{theorem}[Cayley-Hamilton定理]\label{thm:Cayley-Hamilton定理}
    设$\bm{A}\in M_n\left(
        \mathbb{K}
    \right)$,$f\left(\lambda\right)=\left|
    \lambda\bm{I}_n-\bm{A}
    \right|$为其特征多项式,则$f\left(\bm{A}\right)=\bm{O}$.
\end{theorem}
\begin{Proof}
矩阵运算是否为零矩阵与域无关,于是扩张到$\mathbb{C}.$因为任一方阵均复相似于一个上三角阵,于是存在非异阵$\bm{P}\,\textup{s.t.}$
\[
\bm{B}=\bm{P}^{-1}\bm{AP}=\begin{pmatrix}
    \lambda_1&a_{12}&\cdots&a_{1n}\\
    &\lambda_2&\cdots&a_{2n}\\
    &&\ddots&\vdots\\
    &&&\lambda_n
\end{pmatrix}
\]则$f\left(\lambda\right)=
\left(
    \lambda-\lambda_1
\right)\left(
    \lambda-\lambda_2
\right)\cdots\left(
    \lambda-\lambda_n
\right)
    $由上知$f\left(\bm{B}\right)=\bm{O}$,考虑相似关系与多项式相容,得证.
\end{Proof}
\end{formal}
\begin{green}
\begin{corollary}[特征多项式与极小多项式]\label{cor:特征多项式与极小多项式}
    设$\bm{A}\in M_n\left(\mathbb{K}\right)$的特征多项式是$f\left(\lambda\right)$,极小多项式是$m\left(\lambda\right)$,则
    \begin{enumerate}[label = {\textup{(\arabic*)}}]
        \item $m\left(\lambda\right)\mid f\left(x\right)$,特别地,$\deg\,m\left(\lambda\right)\leqslant n$
        \item $f\left(\lambda\right),m\left(\lambda\right)$有相同的根（不计重数）
        \item $f\left(\lambda\right)\mid m\left(\lambda\right)^n$
    \end{enumerate}
\end{corollary}
\begin{Proof}
\begin{enumerate}[label = {\textup{(\arabic*)}}]
    \item 由\textup{Cayley-Hamilton}定理知矩阵适合自己的特征多项式,由极小多项式的定义得证.
    \item 因为特征值均为极小多项式的根,结合上一推论得证.
    \item 设特征多项式$f\left(\lambda\right)=\left(\lambda-\lambda_1\right)^{m_1}
     \left(\lambda-\lambda_2\right)^{m_2}\cdots\left(\lambda-\lambda_k\right)^{m_k}
    $其中$\lambda_1,\lambda_2,\cdots,\lambda_k$是$\bm{A}$的全体不同特征值,$m_1,m_2,\cdots,m_k$是对应的代数重数.于是极小多项式$m\left(\lambda\right)=\left(
        \lambda-\lambda_1
    \right)^{r_1}
    \left(
        \lambda-\lambda_2
    \right)^{r_2}\cdots\left(
        \lambda-\lambda_k
    \right)^{r_k}
    $其中$r_1,r_2,\cdots,r_k\in\mathbb{Z}^+.$注意到$m_1+m_2+\cdots+m_k=n\Longrightarrow
    m_i\leqslant n\leqslant n\cdot r_i,\forall 1\leqslant i\leqslant k.$于是$f\left(\lambda\right)\mid m\left(\lambda\right)^n$.\qedhere
\end{enumerate}
\end{Proof}
\end{green}
\begin{brown}
\begin{example}
    设矩阵$\bm{A}$有$n$个不同的特征值$\lambda_1,\lambda_2,\cdots,\lambda_n$.
    
    一方面,$f\left(\lambda\right)
    =\left(
        \lambda-\lambda_1
    \right)\left(
        \lambda-\lambda_2
    \right)\cdots\left(
        \lambda-\lambda_n
    \right)$.
    
    另一方面,$m\left(\lambda\right)=\left( 
        \lambda-\lambda_1
    \right)\left(
        \lambda-\lambda_2
    \right)\cdots\left(
        \lambda-\lambda_n
    \right)$,故$f\left(\lambda\right)=m\left(\lambda\right)$.
\end{example}
\end{brown}
\begin{brown}
\begin{example}
    设$\bm{A}=c\bm{I}_n$,显然$f\left(\lambda\right)=m\left(\lambda\right)^n.$
\end{example}
\end{brown}
\begin{green}
\begin{corollary}[几何意义的Cayley-Hamilton定理]\label{cor:几何意义的Cayley-Hamilton定理}
    设$\bm{\varphi}\in \mathcal{L}\left(V_{\mathbb{K}}^n\right)$,其特征多项式$f\left(\lambda\right)=\left|
         \lambda\bm{I}_V-\bm{\varphi}
    \right|$,则$f\left(\bm{\varphi}\right)=\bm{0}.$
\end{corollary}
\end{green}
我们曾提到过,不可对角化的矩阵远远没有可对角化的矩阵来得多,并且可以用可对角化的矩阵取去逼近不可对角化的矩阵.
\begin{brown}
\begin{example}
    设$\bm{A}=\begin{pmatrix}
        a_{11}  &a_{12}\\
        a_{21}  &a_{22}
    \end{pmatrix}\in M_2\left(\mathbb{C}\right)$,若$\bm{A}$不可对角化则$\lambda_1=\lambda_2.$其特征多项式
    \[
    \left|
    \lambda\bm{I}_2-\bm{A}   
    \right|=\lambda^2-\left(
        a_{11}+a_{22}
    \right)\lambda+\left(
        a_{11}a_{22}-a_{12}a_{21}
    \right)
    \]于是$\varDelta
     = \left(
        a_{11}+a_{22}
    \right)^2-4\left(
        a_{11}a_{22}-a_{12}a_{21}
    \right)=0
    $.不妨记作$f\left(
        a_{11}, a_{12}, a_{21}, a_{22}
    \right).$
    
    考虑映射
    \begin{align*}
        \bm{\varphi}:M_2\left(\mathbb{C}\right)&\longrightarrow \mathbb{C}^4\\
        \begin{pmatrix}
            a_{11}&a_{12}\\
            a_{21}&a_{22}
        \end{pmatrix}&\longmapsto \begin{pmatrix}
            a_{11}\\
            a_{12}\\
            a_{21}\\
            a_{22}
        \end{pmatrix}
    \end{align*}显然线性同构.

    也就是说,若上述矩阵$\bm{A}$不可对角化,其对应的$\mathbb{C}^4$中的点$\left(
        a_{11},a_{12},a_{21},a_{22}
    \right)$必定是$f$的零点即$f\left(
        a_{11},a_{12},a_{21},a_{22}
    \right)=0.$

    从拓扑的角度看,不可对角化的点落在某一个超平面内,我们可以从该平面之外的某一个点（其一定邻域内的点均是可对角化的）逼近该点.
\end{example}
\end{brown}
基于上述分析,可以给出\cref{thm:Cayley-Hamilton定理}Cayley-Hamilton定理的另一种证明.
\begin{brown}
\begin{Proof}
先证可对角化的矩阵,设
\[
\bm{\varLambda}=
\bm{P}^{-1}\bm{AP}=\begin{pmatrix}
    \lambda_1&&&\\
    &\lambda_2&&\\
    &&\ddots&\\
    &&&\lambda_n  
\end{pmatrix}
\]$f\left(\lambda\right)=
\left(
    \lambda-\lambda_1
\right)\left(
    \lambda-\lambda_2
\right)\cdots\left(
    \lambda-\lambda_n
\right)
$于是
\begin{align*}
&f\left(\bm{\varLambda}\right)\\
=&
\begin{pmatrix}
    0&&&\\
    &\lambda_2-\lambda_1&&\\
    &&\ddots&\\
    &&&\lambda_n-\lambda_1
\end{pmatrix}\begin{pmatrix}
    \lambda_1-\lambda_2&&&\\
    &0&&\\
    &&\ddots&\\
    &&&\lambda_n-\lambda_2
\end{pmatrix}
\begin{pmatrix}
    \lambda_1-\lambda_n&&&\\
    &\lambda_2-\lambda_n&&\\
    &&\ddots&\\
    &&&0
\end{pmatrix}\\
=&\bm{O}
\end{align*}于是$f\left(\bm{A}\right)=\bm{O}.$

考虑任意性,设$\bm{A}\in M_n\left(\mathbb{K}\right)$,设非异阵$\bm{P}$使得
\[
\bm{P}^{-1}\bm{AP}=\begin{pmatrix}
    \lambda_1&*&\cdots&*\\
    &\lambda_2&\cdots&*\\
    &&\ddots&\vdots\\
    &&&\lambda_n
\end{pmatrix}
\]$\exists c_1,c_2,\cdots,c_n\in\mathbb{K}\,\textup{s.t.}\forall 0<t\ll 1,\lambda_1+c_1t,
\lambda_2+c_2t,\cdots,\lambda_n+c_nt$
互不相同.然后构造
\[
\bm{A}_t=   
\bm{P}\begin{pmatrix}
    \lambda_1+c_1t&*&\cdots&*\\
    &\lambda_2+c_2t&\cdots&*\\
    &&\ddots&\vdots\\
    &&&\lambda_n+c_nt
\end{pmatrix}\bm{P}^{-1}
\]当$0<t\ll 1$时其特征值互不相同则必定可对角化,特别地,$t=0$时$\bm{A}_0=\bm{A}.$考虑其特征多项式
\[
\left(
    \bm{A}_t-\left(\lambda_1+c_1t\right)\bm{I}_n
\right)\cdots\left(
    \bm{A}_t-\left(\lambda_n+c_nt\right)\bm{I}_n    
\right)=\bm{O},\forall 0<t\ll 1
\]
考虑到连续性,令$t\to 0$则
\[
\left(
    \bm{A}-\lambda_1\bm{I}_n
\right)\cdots\left(
    \bm{A}-\lambda_n\bm{I}_n
\right)=\bm{O}
\]得证.
\end{Proof}
\end{brown}
\newpage
\section{特征值的估计}
\begin{red}
\begin{remark}
    首先说明记号,定义$\displaystyle
    \bm{A}=\left(
        a_{ij}
    \right)_{n\times n}\in M_n\left(\mathbb{C}\right),R_i=\sum_{j\neq i}\left|a_{ij}\right|=\left|
    a_{i1}
    \right|+
    \left|
    a_{i2}
    \right|+\cdots+\left|
    a_{i,i-1}
    \right|+\left|
    a_{i,i+1}
    \right|+
    \cdots+\left|
    a_{in}
    \right|
    $
\end{remark}
\end{red}
\subsection{Gerschgorin圆盘第一定理}
下面给出Gerschgorin圆盘第一定理（戈尔斯哥利戈氏圆盘第一定理,戈氏第一圆盘定理,圆盘定理）.
\begin{formal}
\begin{theorem}[Gerschgorin圆盘第一定理]\label{thm:Gerschgorin圆盘第一定理}
    设$\bm{A}=\left(
        a_{ij}
    \right)_{n\times n}\in M_n\left(\mathbb{C}\right)$,则$\bm{A}$的所有特征值一定落在以下$n$个$\textup{Gerschgorin}$圆盘中
    \[
    \left|
        z-  a_{ii}
    \right|\leqslant R_i,\forall 1\leqslant i\leqslant n
    \]
\end{theorem}
\begin{Proof}
任取$\bm{A}$的特征值$\lambda_0$,其特征向量$\bm{\xi}=\left(
    \xi_1,\xi_2,\cdots,\xi_n
\right)'\neq \bm{0}$即$\bm{A\xi}=\lambda_0\bm{\xi}.$于是有
\[
\begin{cases*}
a_{11}x_1+a_{12}x_2+\cdots+a_{1n}x_n=\lambda_0x_1\\
a_{21}x_1+a_{22}x_2+\cdots+a_{2n}x_n=\lambda_0x_2\\
\qquad\qquad\cdots\cdots\cdots\cdots\\
a_{n1}x_1+a_{n2}x_2+\cdots+a_{nn}x_n=\lambda_0x_n
\end{cases*}
\]因为$\bm{\xi}$的分量重至少有一个不为零,于是考虑设$0<\left|x_r\right|=\max\left\{
    \left|x_1\right|,\left|x_2\right|,\cdots,\left|x_n\right|
\right\}.$考虑
\[
a_{r1}x_1+\cdots+a_{rr}x_r+\cdots+a_{rn}x_n=\lambda_0x_r
\]得到
\begin{align*}
    \left|
        \lambda_0-a_{rr}
    \right|
    \left|x_r\right|&=\left|
        a_{r1}x_1+\cdots+a_{r,r-1}x_{r-1}+a_{r,r+1}x_{r+1}+\cdots+a_{rn}x_n
    \right|\\
    &\leqslant \left|
        a_{r1}x_1
    \right|+\cdots+\left|
        a_{r,r-1}x_{r-1}
    \right|+\left|
        a_{r,r+1}x_{r+1}
    \right|+\cdots+\left|
        a_{rn}x_n
    \right|\\
    &\leqslant R_r\left|x_r\right|
\end{align*}于是$\left|
    \lambda_0-a_{rr}
\right|\leqslant R_r$.
\end{Proof}
\end{formal}
\begin{brown}
\begin{example}
    估计矩阵的特征值
    \[
    \begin{pmatrix}
        1&0.5&-0.2&-1\\
        0.3&2&-0.2&1.1\\
        -0.5&0.1&-4&0.2\\
        -1&-0.1&0.2&0
    \end{pmatrix}
    \]
\end{example}
\begin{solution}
    明显地
    \begin{align*}
        &D_1: \left|
            z-1
        \right|\leqslant 0.5+0.2+1=1.7\\
        &D_2: \left|
            z-2
        \right|\leqslant 0.3+0.2+1.1=1.6\\
        &D_3: \left|
            z+4
        \right|\leqslant 0.5+0.1+0.2=0.8\\
        &D_4: \left|
            z
        \right|\leqslant 1+0.1+0.2=1.3
    \end{align*}
\end{solution}
\end{brown}
\begin{formal}
\begin{definition}[连通区域]\label{def:连通区域}
    若几个戈氏圆盘相连在一起,则称相连区域（指它们覆盖的所有区域）为连通区域,这几个圆盘为连通圆盘.
\end{definition}
\end{formal}
\begin{red}
\begin{remark}
    对于单个非连通的圆盘,该圆盘内必定存在一个特征值.但对于连通区域中的某一圆盘,只能保证特征值在该连通区域内但不一定在这个圆盘内.这将在\textup{Gerschgorin}圆盘第二定理中体现.
\end{remark}
\end{red}
\subsection{Gerschgorin圆盘第二定理}
\begin{formal}
\begin{theorem}[多项式系数]\label{thm:多项式系数}
    设$f\left(x\right)=a_nx^n+
    a_{n-1}x^{n-1}+\cdots+a_1x+a_0
    $,则$f\left(x\right)$的根关于$a_n,\cdots,a_1,a_0$连续.
\end{theorem}
\end{formal}
\begin{red}
\begin{remark}
    我们曾说过,对于高于$4$次的多项式,没有能够用系数的初等运算得到的显式求根公式,也无从谈起像二次多项式那样$\displaystyle
    \frac{
        -b\pm\sqrt{b^2-4ac}
    }{2a}$的连续性,因此这里需要做一个解释.

    设$f\left(x\right)$的根为$\lambda_1,\lambda_2,\cdots,\lambda_n$,则$\lambda_i=\lambda_i\left(
        a_n,\cdots,a_1,a_0
    \right),\forall 
    1\leqslant i\leqslant n
    $,考虑
    \[
    \widetilde{f}\left(x\right)=\widetilde{a}_nx^n+\widetilde{a}_{n-1}x^{n-1}+\cdots+\widetilde{a}_1x+\widetilde{a}_0
    \]的根为
    $\widetilde{\lambda}_1,\widetilde{\lambda}_2,\cdots,\widetilde{\lambda}_n$,则$\widetilde{\lambda}_i=\widetilde{\lambda}_i\left(
        \widetilde{a}_n,\cdots,\widetilde{a}_1,\widetilde{a}_0
    \right),\forall
    1\leqslant i\leqslant n$.
    
    于是$\forall \varepsilon >0,\exists \delta$使得当$\left|
        \widetilde{a}_n-a_n
    \right|<\delta
    ,\cdots,\left|
        \widetilde{a}_1-a_1
    \right|<\delta,\left|
        \widetilde{a}_0-a_0
    \right|<\delta$时
    ,必定有$\left|
        \widetilde{\lambda}_i\left(
            \widetilde{a}_n,\cdots,\widetilde{a}_1,\widetilde{a}_0
        \right)-\lambda_i\left(
            a_n,\cdots,a_1,a_0
        \right)
    \right|<\varepsilon,\forall 1\leqslant i\leqslant n$.
    
    上定理证明需要利用复变函数论中的\textup{Rouch\'e}定理,于是这里不再赘述.可参考\textup{\cite{蒋尔雄1978}}中的证明.
\end{remark}
\end{red}
\begin{red}
\begin{remark}
    上述定理并非断言多项式的某一个根是关于系数的连续函数,因此一般而言,矩阵的特征值也并不是关于矩阵元素的连续函数.
    
    例如$\bm{A}=\begin{pmatrix}
        0&z\\1&0
    \end{pmatrix}$,当$z$在复平面上的单位圆内变化时,其特征值不是关于$z$的连续函数,但对于$z$在实轴上的某一区间来说,是成立的,这证明下面\textup{Gerschgorin}圆盘第二定理利用摄动法证明所用到的.
\end{remark}
\end{red}
下面给出Gerschgorin圆盘第二定理,该定理能够对连通区域内特征值作进一步的估计.
\begin{red}
    \begin{remark}
        戈氏圆盘可能是相同的,也可能是不同的.
        相同时在下面的第二定理中需要计算重数.
    \end{remark}
\end{red}
\begin{formal}
\begin{theorem}[Gerschgorin圆盘第二定理]\label{thm:Gerschgorin圆盘第二定理}
    设$k$个戈氏圆盘连成一个连通区域,则有且仅有$k$个特征值落在该连通区域内.

    特别地,重合的圆盘计重数,相同的特征值亦计重数.
\end{theorem}
\begin{Proof}
考虑摄动法,设$\bm{A}=\left(a_{ij}\right)_{n\times n}$,考虑摄动
\[
\bm{A}\left(0\right)=\begin{pmatrix}
    a_{11}&&&\\
    &a_{22}&&\\
    &&\ddots&\\
    &&&a_{nn}
\end{pmatrix}
\]而
\[
\bm{A}\left(t\right)=\begin{pmatrix}
    a_{11}&ta_{12}&\cdots&ta_{1n}\\
    ta_{21}&a_{22}&\cdots&ta_{2n}\\
    \vdots&\vdots&\ddots&\vdots\\
    ta_{n1}&ta_{n2}&\cdots&a_{nn}
\end{pmatrix}
\]$\bm{A}\left(1\right)=\bm{A}.$根据\textup{Gerschgorin}圆盘第一定理,$\bm{A}\left(t\right)$的特征值落在
\[
\left|
    z-a_{ii}
\right|\leqslant t\cdot R_i,\forall 
1\leqslant i\leqslant n
\]内.因此$\forall 0\leqslant t\leqslant 1,\bm{A}\left(t\right)$的特征值均落在$\bm{A}$的对应圆盘中
\[
\left|
    z-a_{ii}  
\right|\leqslant
R_i,\forall 1\leqslant i\leqslant n
\]考虑$\bm{A}\left(t\right)$的特征多项式,其系数均是$t$的多项式,于是由\cref{thm:多项式系数}知其关于$t$连续.于是特征值关于系数连续,系数关于$t$连续,因此有$\bm{A}\left(t\right)$的特征值关于$t$连续.

断言,特征值在连续变化的过程中,只能在同一个连通区域内变化而不能跳出该连通区域.考虑反证法,如果变化过程中跳出,意味着$\exists t_0\in \left(0,1\right)$使得$\lambda_i\left(t_0\right)$不落在任何一个戈氏圆盘中.这与上文$\forall 0\leqslant t\leqslant 1,\bm{A}\left(t\right)$的特征值均落在$\bm{A}$的对应圆盘中
\[
\left|
    z-a_{ii}  
\right|\leqslant
R_i,\forall 1\leqslant i\leqslant n
\]矛盾.

于是从$t:0\to 1$,直到$t=1$,对于一个有$k$个圆盘连通的区域,其内至少有$t=0$时的$k$个特征值,另一方面,该断言对于所有连通区域均成立,于是也至多是$k$个,故有且仅有$k$个特征值落在该连通区域内.证毕.
\end{Proof}
\end{formal}
\begin{brown}
\begin{example}
    考虑矩阵$\bm{A}=\left(a_{ij}\right)_{n\times n}$且$0\ll a_{11}\ll a_{22}\ll\cdots\ll a_{nn}$,则$\bm{A}$可对角化.
\end{example}
\end{brown}
\newpage
\section{总结}
\begin{brown}
\begin{example}
    设$\bm{A}^2-\bm{A}-3\bm{I}_n=\bm{O}$,求证：$\bm{A}-2\bm{I}_n$可逆.
\end{example}
\begin{Proof}
\begin{enumerate}[label={\textup{(
    \arabic*
)}}]
    \item 凑因子：\[
    \left(\bm{A}-2\bm{I}_n\right)\left(\bm{A}+\bm{I}_n\right)=\bm{I}_n
    \]
    \item 求解方程组：只要证$\left(
        \bm{A}-2\bm{I}_n
    \right)\bm{x}=\bm{0}$只有零解
    \item 互素多项式：\[
    \left(
        x^2-x-3,x-2
    \right)=1
    \]
    \item 特征值：反证法,考虑设$\bm{A}-2\bm{I}_n$不可逆,即其行列式为零,则$2$是$\bm{A}$的一个特征值,但是该特征值却不适合多项式$x^2-x-3$,矛盾.\qedhere
\end{enumerate}
\end{Proof}
\end{brown}
\begin{brown}
\begin{example}
    设$4$阶方阵$\bm{A}$满足$\mathrm{tr}\,\bm{A}^i=i,i=1,2,3,4$.求$\left|\bm{A}\right|$.
\end{example}
\begin{solution}
    设$\bm{A}$的特征值$\lambda_1,\lambda_2,\lambda_3,\lambda_4$,则
    $\bm{A}^i$的特征值为$\lambda_1^i,\lambda_2^i,\lambda_3^i,\lambda_4^i$,于是$\displaystyle
    \sum_{k=1}^{4}\lambda_k^i=i$即\textup{Fermal}和$s_1=1,s_2=2,s_3=3,s_4=4.$

    因为$\left|\bm{A}\right|=\lambda_1\lambda_2\lambda_3\lambda_4=\sigma_4$.于是由牛顿公式,$k\leqslant n=4$时
    \[
    s_k-s_{k-1}\sigma_1+\cdots+\left(-1\right)^kk\sigma_k=0
    \]
    于是取$k=1,2,3,4$得到$\displaystyle
    \sigma_1=1,\sigma_2=-\frac{1}{2},\sigma_3=\frac{1}{6},\sigma_4=\frac{1}{24}.$
\end{solution}
\end{brown}
\begin{brown}
\begin{example}
    设矩阵$\bm{A}^m,\bm{B}^n$且无公共特征值,证明矩阵方程$\bm{AX}=\bm{XB}$只有零解.
\end{example}
\begin{Proof}
    利用\cref{thm:Cayley-Hamilton定理}\textup{Cayley-Hamilton}定理,设$\bm{A}$的特征多项式为$f\left(\lambda\right)=\left|\lambda\bm{I}_m-\bm{A}\right|$,于是$f\left(\bm{A}\right)=\bm{O}$.同时容易有$\bm{A}^2\bm{X}=\bm{XB}^2,\cdots$于是有$\bm{O}=f\left(\bm{A}\right)\bm{X}=\bm{X}f\left(\bm{B}\right).$
    
    下面考虑证明$f\left(\bm{B}\right)$非异,不妨设$\bm{B}$的特征值$\mu_1,\mu_2,\cdots,\mu_n$则$f\left(\bm{B}\right)$的特征值为$f\left(\mu_1\right)\neq 0,f\left(\mu_2\right)\neq 0,\cdots,f\left(\mu_n\right)\neq 0$,那么$f\left(\bm{B}\right)$可逆.于是$\bm{X}=\bm{O}$得证.

    也可设$\bm{B}$特征多项式$g\left(\lambda\right)$且$\left(
        f\left(\lambda\right),g\left(\lambda\right)
    \right)=1$,考虑带余除法则
    \[
    f\left(\lambda\right)u\left(\lambda\right)+g\left(\lambda\right)v\left(\lambda\right)=1
    \]于是$f\left(\bm{B}\right)u\left(\bm{B}\right)=\bm{I}_n$得证$f\left(\bm{B}\right)$可逆.
\end{Proof}
\end{brown}
\begin{brown}
\begin{example}
    设$n$阶方阵$\bm{A},\bm{B}$,特征值均大于零,且$\bm{A}^2=\bm{B}^2.$求证：$\bm{A}=\bm{B}$.
\end{example}
\begin{Proof}
    即上例中$\bm{X}=\bm{A}-\bm{B}$,因为
    \[
    \bm{A}\left(\bm{A}-\bm{B}\right)=
    \left(\bm{A}-\bm{B}\right)\left(-\bm{B}\right)
    \]考虑到$\bm{A}$和$-\bm{B}$没有公共特征值,由上例得证.
\end{Proof}
\end{brown}
\begin{brown}
\begin{example}
    设$n$阶方阵$\bm{A}$的特征值均为偶数,证明：$\bm{X}+\bm{AX}=\bm{XA}^2$只有零解.
\end{example}
\begin{Proof}
显然考虑
\[
\left(\bm{A}+\bm{I}\right)\bm{X}=\bm{X}\bm{A}^2
\]由上例,考虑$\bm{A}+\bm{I}$与$\bm{A}^2$的特征值即可,显然成立.
\end{Proof}
\end{brown}
\begin{brown}
\begin{example}
    设$n$阶方阵$\bm{A}$适合$a_mx^m+
    \cdots+a_1x+a_0
    $且$\displaystyle
    \left|a_m\right|\geqslant
    \sum_{i=0}^{m-1}\left|a_i\right|
    $.证明：$2\bm{X}+\bm{AX}=\bm{XA}$仅有零解.
\end{example}
\begin{Proof}
任取一特征值$\lambda_0$,断言$\left|\lambda_0\right|<1$,反证法,设$\left|\lambda_0\right|\geqslant 1.$因为$\lambda_0$适合该多项式即
\[
a_m\lambda_0^m+\cdots+a_1\lambda_0+a_0=0
\]
于是
\[
a_m=-\frac{a_{m-1}}{\lambda_0}-\cdots-\frac{a_1}{\lambda_0^{m-1}}-\frac{a_0}{\lambda_0^m}
\]
两边取模长\newpage
\[
\left|a_m\right|\leqslant
\sum_{i=0}^{m-1}\left|a_i\right|
\]显然矛盾.

于是来考虑$\bm{A}$和$\bm{A}+2\bm{I}_n$的特征值.因为$\bm{A}$的特征值均在复平面单位圆内部,而$\bm{A}+2\bm{I}_2$的特征值在向右平移两个单位的圆内且均不包含边界于是二者没有公共特征值.于是得证.
\end{Proof}
\end{brown}
\begin{brown}
\begin{example}
    设矩阵$\bm{A}^m\in M_m\left(
        \mathbb{K}
    \right),\bm{B}^n\in M_n\left(\mathbb{K}\right)$无公共特征值,则$\forall \bm{C}\in M_{m\times n}\left(\mathbb{K}\right)$,矩阵方程$\bm{AX}-\bm{XB}=\bm{C}$有唯一解.
\end{example}
\begin{Proof}
考虑$\bm{\varphi}\in \mathcal{L}\left(
    M_{m\times n}\left(\mathbb{K}\right)
\right):\bm{X}\longmapsto
\bm{AX}-\bm{XB}$,由上例可知$\bm{\varphi}$是单射,那么也是满射,于是是线性同构.于是$\forall \bm{C}\in M_{m\times n}\left(\mathbb{K}\right)$,一定存在唯一一个$\bm{X}_0\in M_{m\times n}\left(\mathbb{K}\right)$使得$\bm{\varphi}\left(\bm{X}_0\right)=\bm{AX}_0-\bm{X}_0\bm{B}=\bm{C}.$证毕.
\end{Proof}
\end{brown}
\begin{brown}
\begin{example}
    设矩阵$\bm{A}^m,\bm{B}^n$无公共特征值且均可对角化,则$\bm{M}=\begin{pmatrix}
        \bm{A}&\bm{C}\\\bm{O}&\bm{B}
    \end{pmatrix}$可对角化.
\end{example}
\begin{Proof}
由上例知存在$\bm{X}_0$使得$\bm{AX}_0-\bm{X}_0\bm{B}=\bm{C}$.于是
\begin{align*}
    \begin{pmatrix}
        \bm{I}_m&\bm{X}_0\\\bm{O}&\bm{I}_n
    \end{pmatrix}\begin{pmatrix}
        \bm{A}&\bm{C}\\\bm{O}&\bm{B}
    \end{pmatrix}\begin{pmatrix}
        \bm{I}_m&-\bm{X}_0\\\bm{O}&\bm{I}_n
    \end{pmatrix}=\begin{pmatrix}
        \bm{A}&\bm{O}\\\bm{O}&\bm{B}
    \end{pmatrix}
\end{align*}
相似得证,于是$\bm{M}$可对角化.
\end{Proof}
\end{brown}
\begin{brown}
\begin{example}
    设$\bm{A}$为非异循环矩阵,求证：$\bm{A}^{-1}$也是循环矩阵.
\end{example}
\begin{Proof}
法一,设\newpage\[
\bm{A}=\begin{pmatrix}
    a_1&a_2&\cdots&a_n\\
    a_n&a_1&\cdots&a_{n-1}\\
    \vdots&\vdots&&\vdots\\
    a_2&a_3&\cdots&a_1
\end{pmatrix}
\]并有基础循环矩阵$\bm{J}=\begin{pmatrix}
    &\bm{I}_{n-1}\\1&
\end{pmatrix}$则$\bm{A}=a_1\bm{I}_n+a_2\bm{J}+\cdots+a_n\bm{J}^{n-1}$.作$g\left(x\right)=a_1+a_2x+\cdots+a_nx^{n-1}$且$\bm{J}$的特征多项式$\displaystyle
f\left(\lambda\right)=\lambda^n-1\Longrightarrow\omega_k=e^{\frac{2k\pi \mathrm{i}}{n}},0\leqslant k\leqslant n-1
$.因此$\bm{A}$的特征值是$g\left(\omega_k\right),0\leqslant k\leqslant n-1$即特征值,因为$\bm{A}$非异即$g\left(\omega_k\right)$非零即推知$\left(
    f\left(x\right),g\left(x\right)
\right)=1$即带余除法
\[
f\left(x\right)u\left(x\right)+g\left(x\right)v\left(x\right)=1
\]代入$\bm{J}$即得$\bm{A}^{-1}=v\left(\bm{J}\right).$

法二,因为$\bm{A}$可逆则$\bm{A}^{-1}=h\left(\bm{A}\right)=h\left(g\left(\bm{J}\right)\right)$得证.
\end{Proof}
\end{brown}